% !TEX root = ../../../main.tex
% Block card: detection.mp_motifs -- "Motif groups" (Part I, the blocks chapter).
\cardfield{What it does}
Reads a matrix profile and lists the shapes that repeat best: for each, the window where it
occurs most cleanly and that window's few closest look-alikes. Each such set is a group, a
candidate \term{motif}{motif}.

\cardfield{When to use it}
Straight after \blockref{detection.matrix_profile}, when you want the repeated shapes themselves
rather than a curve of distances. It is the matrix-profile route to motifs; if you already have an
example of the shape you want, \blockref{detection.seed_matches} searches for it directly.

\cardfield{What goes in, what comes out}
In goes the matrix profile. Out come stretches of the recording, each one window long, labelled
by group: \texttt{G0:seed} is the best-matched window of the most repeated shape, \texttt{G0:nn1}
and \texttt{G0:nn2} its closest look-alikes, then \texttt{G1:\dots} for the next shape, and so on.
Each stretch carries its distance from the group's first window. They are saved as detections.

\cardfield{Settings that matter}
\emph{max\_motifs} is how many groups to list, best first. \emph{n\_neighbors} is how many matches
to fetch per group, \emph{counting the starting window itself}: the default 3 gives the start and
at most two look-alikes. \emph{max\_distance} drops look-alikes that are too far to count (0 keeps
them all).

\cardfield{Watch out for}
Two look-alikes in one group can overlap by up to three quarters of a window: often the same event,
shifted a little, rather than two events. Different groups can also share windows, so one event
can appear in two groups. A window marks where the match is best, which may start part-way into
the event. And if the recording was not detrended first, the best ``shapes'' are usually drift.
